% Chapter 8, Section 2

\section{Momentum-Based Methods}
\label{sec:momentum}

\subsection{Intuition: Rolling a Ball Down a Valley}

Plain SGD can wobble like a marble on a bumpy path. \textbf{Momentum} acts like mass: it carries velocity so you keep moving in consistently good directions and smooth out small bumps. \textbf{Nesterov acceleration} adds anticipation by peeking where the momentum will take you before correcting, often yielding crisper steps in curved valleys.\index{momentum}\index{Nesterov}

Historical note: Momentum has deep roots in convex optimization and was popularized in early neural network training; Nesterov's variant provided stronger theoretical guarantees in convex settings and inspired practical variants in deep learning \cite{Polyak1964,Nesterov1983,GoodfellowEtAl2016,Bishop2006}.

\subsection{Momentum}

Accumulates gradients over time:
\begin{align}
\vect{v}_t &= \beta \vect{v}_{t-1} - \alpha \nabla_{\vect{\theta}} L(\vect{\theta}_t) \\
\vect{\theta}_{t+1} &= \vect{\theta}_t + \vect{v}_t
\end{align}

where $\beta \in [0, 1)$ is the momentum coefficient (typically 0.9).

Benefits:\index{optimization!momentum}
\begin{itemize}
    \item Accelerates convergence in relevant directions
    \item Dampens oscillations
    \item Helps escape local minima and saddle points
\end{itemize}

\subsection{Nesterov Accelerated Gradient (NAG)}\index{optimization!Nesterov}

"Look-ahead" version of momentum:
\begin{align}
\vect{v}_t &= \beta \vect{v}_{t-1} - \alpha \nabla_{\vect{\theta}} L(\vect{\theta}_t + \beta \vect{v}_{t-1}) \\
\vect{\theta}_{t+1} &= \vect{\theta}_t + \vect{v}_t
\end{align}

Evaluates gradient at anticipated future position, often providing better updates.

